\section{Problem Statement}
\label{sec:problem}

We study the following \emph{function-level fault localisation} task: given
(i)~an entry function from which the bug manifests, (ii)~a natural-language
description of the observed symptom, and (iii)~optional access to a call-graph
of the target codebase, identify the single function most likely to be faulty.

\paragraph{Target program.}
We use \sqlite~3.30.1 (commit \texttt{25e9591}, tag \texttt{version-3.30.1}),
a widely studied embedded SQL engine~\cite{sqlite}.
We extract call graphs from eleven core C source files spanning tokenization,
parsing, query planning, expression resolution, and virtual machine execution
(348{,}852 tokens, approximately 33{,}000 lines of C).
\sqlite is a suitable subject because its source is publicly available, it is
used in real-world FL research, and its function-pointer-heavy dispatch tables
directly expose the structural-gap challenge studied here.

\paragraph{Benchmark.}
We evaluate on six questions drawn from the course benchmark, each specifying
an entry function, a symptom description, and a ground-truth function.
The targets span SQL query processing (\texttt{sqlite3Select},
\texttt{flattenSubquery}, \texttt{withExpand}), string formatting
(\texttt{sqlite3\_str\_vappendf}), name resolution
(\texttt{resolveOrderByTermToExprList}), and aggregate accumulation
(\texttt{resetAccumulator}).
Two questions (Q1 and Q4) target functions that are reachable from the entry
only via function-pointer calls and are therefore absent from both graphs.

\paragraph{Structural challenge.}
\sqlite relies heavily on function pointers stored in virtual-table and callback
structs.
Neither the compile-time nor the syntactic graph models these dynamic edges, creating
structural gaps between the entry function and certain targets.
While a pure graph traversal alone along static edges cannot reach these targets, an agent
equipped with source-reading and symbol-search tools can in principle bridge the gap if
given sufficient search budget.
